\section{Ciclos de vida}
\label{sec:estados}

\subsection{Reserva}
\begin{itemize}
    \item Pendiente: solicitud registrada sin bloqueo de disponibilidad.
    \item Confirmada: solicitud aprobada; el intervalo queda comprometido para la unidad y el solicitante.
    \item Rechazada: solicitud pendiente denegada con motivo.
    \item Cancelada: solicitud pendiente o confirmada retirada antes de la entrega.
    \item Atendida: reserva confirmada convertida en su \'unico pr\'estamo.
    \item Vencida: reserva pendiente o confirmada cuyo intervalo termin\'o sin entrega.
\end{itemize}

Las transiciones permitidas son pendiente a confirmada, rechazada, cancelada o vencida, y confirmada a atendida, cancelada o vencida. Los estados terminales no se reabren; una nueva solicitud genera otra reserva. El paso a vencida debe efectuarse antes de evaluar operaciones que dependan de su vigencia, aunque no exista un proceso peri\'odico en ejecuci\'on.

\subsection{Pr\'estamo}
\begin{itemize}
    \item Activo: entrega registrada, sin devoluci\'on ni cierre y con plazo a\'un no superado.
    \item Vencido: sigue abierto y la fecha y hora actuales superan el vencimiento vigente.
    \item Devuelto: se registr\'o la recepci\'on e inspecci\'on del bien.
    \item Cerrado por p\'erdida: el administrador document\'o la imposibilidad de devoluci\'on y dio de baja la unidad.
\end{itemize}

Las transiciones son activo a vencido, devuelto o cerrado por p\'erdida, y vencido a devuelto o cerrado por p\'erdida. Una renovaci\'on v\'alida mantiene el estado activo y cambia el vencimiento; no permite reactivar un pr\'estamo ya vencido. El estado vencido debe derivarse del reloj y del cierre, o actualizarse de forma consistente antes de cada consulta u operaci\'on. No puede depender de una tarea manual.

\subsection{Unidad y obligaciones asociadas}
La entrega cambia la unidad de disponible a prestada en la misma transacci\'on que crea el pr\'estamo. Mientras este est\'a activo o vencido, la unidad permanece prestada. La devoluci\'on la lleva a disponible, mantenimiento o baja seg\'un la inspecci\'on; el cierre por p\'erdida la lleva a baja. Una unidad en mantenimiento puede volver a disponible tras verificaci\'on administrativa. Una baja es terminal dentro de esta versi\'on.

Cerrar el pr\'estamo no equivale a liberar toda obligaci\'on: las garant\'ias, incidencias y sanciones conservan sus propios estados hasta su resoluci\'on. Se permite registrar esa resoluci\'on aunque la cuenta del titular est\'e inactiva.

\section{Reglas de integridad}
\label{sec:integridad}

\subsection{Identidad y referencias}
\begin{itemize}
    \item Son \'unicos el documento de persona, nombre de usuario, referencia de persona en la cuenta, c\'odigos de perfiles, nombre de tipo de bien y c\'odigo de inventario.
    \item La referencia a reserva en un pr\'estamo es opcional y \'unica cuando existe; la referencia a pr\'estamo en una garant\'ia es obligatoria y \'unica.
    \item La serie no es obligatoria para todos los bienes. Se controla la combinaci\'on marca, modelo y serie cuando se registra; el ISBN puede compartirse entre ejemplares.
    \item Solicitante, unidad, autorizante, inicio, vencimiento y condiciones aplicadas son obligatorios en todo pr\'estamo.
    \item Las referencias hist\'oricas se conservan mediante desactivaci\'on o baja l\'ogica; no se eliminan en cascada pr\'estamos, garant\'ias, sanciones o eventos.
\end{itemize}

\subsection{Tiempo, disponibilidad y cupo}
Dos intervalos $[i_1,f_1)$ y $[i_2,f_2)$ se solapan si y solo si:
\[
    i_1 < f_2 \quad \land \quad i_2 < f_1.
\]

\begin{itemize}
    \item Toda reserva tiene fin posterior a inicio. Todo pr\'estamo tiene vencimiento posterior a entrega. Devoluci\'on y cierre por p\'erdida, cuando existan, no pueden ser anteriores a la entrega y son mutuamente excluyentes.
    \item Un pr\'estamo devuelto exige fecha de devoluci\'on; uno cerrado por p\'erdida exige fecha de cierre e incidencia. Los abiertos no tienen ninguna de esas fechas.
    \item No se permiten reservas confirmadas solapadas para una unidad. Un pr\'estamo o renovaci\'on no puede invadir otra reserva confirmada; la reserva de origen se excluye al convertirla, para no contar dos veces la misma asignaci\'on.
    \item Un pr\'estamo abierto ocupa su intervalo hasta el vencimiento previsto para la planificaci\'on. Si ya est\'a vencido, la disponibilidad futura se considera bloqueada hasta la devoluci\'on o cierre efectivo. Nunca se entrega otra vez una unidad todav\'ia prestada.
    \item El cupo simult\'aneo se valida sobre todas las unidades del solicitante, no por cada tipo de manera independiente. Para intervalos futuros se suman las reservas confirmadas y los pr\'estamos abiertos que los intersectan; los vencidos se consideran abiertos sin fin conocido.
    \item La confirmaci\'on de reservas, las entregas y las renovaciones verifican tanto conflictos de unidad como de cupo dentro de la misma transacci\'on. Una consulta previa de disponibilidad no reemplaza esta verificaci\'on.
    \item Duraciones y cupos son positivos; tolerancia, tarifa, monto y n\'umero de renovaciones son no negativos. Si no se permiten renovaciones, su m\'aximo es cero.
\end{itemize}

\subsection{Garant\'ias, sanciones y auditor\'ia}
\begin{itemize}
    \item Una entrega que exige garant\'ia no puede confirmarse sin registrarla como vigente. La liberaci\'on exige fecha y responsable; una p\'erdida exige una incidencia vinculada.
    \item El usuario sancionado debe coincidir con el solicitante del pr\'estamo, nunca con el administrador que lo autoriz\'o por el solo hecho de autorizarlo.
    \item El c\'alculo autom\'atico de retraso es idempotente: repetir el procesamiento no duplica sanciones. Emplea la tarifa y tolerancia conservadas, y el vencimiento vigente.
    \item Una sanci\'on no se marca cumplida mientras quede una obligaci\'on econ\'omica o temporal pendiente. Una apelaci\'on y su resoluci\'on no borran el registro original.
    \item Los cambios de estado y sus eventos se confirman juntos. Un fallo no puede dejar un pr\'estamo creado sin ocupar la unidad, una reserva atendida sin pr\'estamo o una garant\'ia sin asociaci\'on.
\end{itemize}

La implementaci\'on deber\'a combinar restricciones declarativas, \'indices y validaciones transaccionales. Una restricci\'on sobre una sola fila no basta para impedir todos los solapamientos ni excesos de cupo.
